% Minimal retained preamble for lane-p2 cards taken from
% sources/logic-fl05-ramsey-like-theorems-for-the-schre/pinned-original-source.tex
% (arXiv:2412.11598v3, CC BY 4.0).  This source ships no separate preamble file:
% the pinned tex is a concatenation of the authors' main.tex and macros.tex and
% carries its preamble inline, so the blocks below are sliced mechanically from
% that file and nothing has been renamed or redefined.
%
% Copied verbatim from the source:
%   line 2        : the authors' class line
%   lines 3-5,7,9 : graphicx, amsthm, the ams bundle, cleveref, url
%   line 10       : tikz-cd, used by the authors' commutative diagrams
%   lines 20-108  : the authors' macro definitions, taken in full so that every
%                   macro of the retained spans is available
%   lines 113-124 : the authors' theorem declarations
%   lines 131-141 : the authors' machinery for restating a theorem
%   lines 143-148 : the authors' coloured comment macros
%
% Omitted source preamble lines and why:
%   line 1, 110-111 and 150 (concatenation comment, title, date, author): the
%       assembler supplies the title block, author list and date;
%   line 6 (hyperref) and line 8 (inputenc with the utf8x option): the forced
%       xelatex engine supplies the hyperref wiring itself and utf8x is
%       incompatible with it;
%   line 115 (numberwithin theorem by section): the assembler numbers the
%       retained material itself;
%   lines 120 and 121-123 are declarations that the authors left commented out or
%       that duplicate earlier ones and are not reproduced.
\documentclass{article}
\usepackage{graphicx}
\usepackage{amsthm}
\usepackage{amsfonts, amsmath, amssymb, amsthm}
\usepackage{cleveref}
\usepackage{url}
\usepackage{tikz-cd}
\newcommand{\Psf}{\mathsf{P}}
\newcommand{\Qsf}{\mathsf{Q}}
\newcommand{\Rsf}{\mathsf{R}}
\def\F{\mathcal{F}}
\def\U{\mathcal{U}}
\def\C{\mathcal{C}}
\def\I{\mathcal{I}}
\def\D{\mathcal{D}}
\def\L{\mathcal{L}}
\def\V{\mathcal{V}}
\def\A{\mathcal{A}}
\def\B{\mathcal{B}}
\def\P{\mathcal{P}}
\def\M{\mathcal{M}}
\def\Nc{\mathcal{N}}
\def\Q{\mathcal{Q}}
\newcommand{\cs}{2^{\NN}}
\newcommand{\bs}{\NN^{\NN}}
\newcommand{\bstr}{2^{<\NN}}
\newcommand{\uh}{\upharpoonright}
\newcommand{\converge}{\!\!\downarrow}
\newcommand{\diverge}{\!\!\uparrow}
\newcommand{\qvdash}{\operatorname{{?}{\vdash}}}
\newcommand{\nqvdash}{\operatorname{{?}{\nvdash}}}
\newcommand{\ISig}{\mathsf{I}\Sigma^0}
\newcommand{\Robinson}{\mathsf{Q}}
\newcommand{\IPi}{\mathsf{I}\Pi^0}
\newcommand{\IDelta}{\mathsf{I}\Delta^0}
\newcommand{\Iopen}{\mathsf{I_{open}}}
\newcommand{\BSig}{\mathsf{B}\Sigma^0}
\newcommand{\BPi}{\mathsf{B}\Pi^0}
\newcommand{\BDelta}{\mathsf{B}\Delta^0}
\newcommand{\BCSig}{\mathsf{BC}\Sigma^0}
\newcommand{\BCPi}{\mathsf{BC}\Pi^0}
\newcommand{\BCDelta}{\mathsf{BC}\Delta^0}
\newcommand{\LSig}{\mathsf{L}\Sigma^0}
\newcommand{\LPi}{\mathsf{L}\Pi^0}
\newcommand{\CSig}{\mathsf{C}\Sigma}
\newcommand{\RCA}[0]{\mathsf{RCA}}
\newcommand{\WKL}[0]{\mathsf{WKL}}
\newcommand{\WWKL}[0]{\mathsf{WWKL}}
\newcommand{\ACA}[0]{\mathsf{ACA}}
\newcommand{\ATR}[0]{\mathsf{ATR}}
\newcommand{\PIOOCA}[0]{\Pi^1_1\mathsf{-CA}}
\newcommand{\CAC}[0]{\mathsf{CAC}}
\newcommand{\ADS}[0]{\mathsf{ADS}}
\newcommand{\RRT}[0]{\mathsf{RRT}}
\newcommand{\FS}[0]{\mathsf{FS}}
\newcommand{\PFS}[0]{\mathsf{PFS}}
\newcommand{\TS}[0]{\mathsf{TS}}
\newcommand{\GRRT}[0]{\mathsf{GRRT}}
\newcommand{\LRRT}[0]{\mathsf{LRRT}}
\newcommand{\EM}[0]{\mathsf{EM}}
\newcommand{\SEM}[0]{\mathsf{SEM}}
\newcommand{\HalfEM}[0]{\mathsf{Half}\mbox{-}\mathsf{EM}}
\newcommand{\RT}[0]{\mathsf{RT}}
\newcommand{\SRT}[0]{\mathsf{SRT}}
\newcommand{\NN}[0]{\mathbb{N}}
\newcommand{\nat}[0]{\mathbb{N}}
\newcommand{\card}{\operatorname{card}}
\newcommand{\finsub}{\subseteq_{\mathtt{fin}}}
\newcommand{\W}{\mathrm{W}}
\newcommand{\sW}{\mathrm{sW}}
\newcommand{\lex}{\mathrm{lex}}
\newcommand{\base}{\operatorname{base}}
\newcommand{\range}{\operatorname{ran}}
\def\qt#1{``#1''}%

\newtheorem*{statement}{Statement}
\newtheorem{theorem}{Theorem}
\numberwithin{theorem}{section}
\newtheorem{maintheorem}[theorem]{Main Theorem}
\newtheorem{lemma}[theorem]{Lemma}
\newtheorem{question}[theorem]{Question}
\newtheorem{proposition}[theorem]{Proposition}
%\newtheorem{statement}[theorem]{Statement}
\newtheorem{remark}[theorem]{Remark}
\newtheorem{definition}[theorem]{Definition}
\newtheorem{corollary}[theorem]{Corollary}
\newtheorem{example}[theorem]{Example}

% Allow to repeat theorems
\makeatletter
\newtheorem*{rep@theorem}{\rep@title}
\newcommand{\newreptheorem}[2]{%
\newenvironment{rep#1}[1]{%
 \def\rep@title{#2 \ref{##1}}%
 \begin{rep@theorem}}%
 {\end{rep@theorem}}}
\makeatother

\newreptheorem{theorem}{Theorem}

% Comments
\usepackage{xcolor}
\newcommand{\ludovic}[1]{\textcolor{blue}{#1}}
\newcommand{\quentin}[1]{\textcolor{green}{#1}}
\newcommand{\lorenzo}[1]{\textcolor{orange}{#1}}
\newcommand{\oriola}[1]{\textcolor{red}{#1}}
\crefname{originalproposition}{original proposition}{original propositions}
\Crefname{originalproposition}{Original proposition}{Original propositions}
\makeatletter
\newlabel{prop:fs-n-k-strong-cone-avoidance}{{6.8}{source}{}{sourcepointer-1}{}}
\newlabel{prop:fs-n-k-strong-cone-avoidance@cref}{{[originalproposition][0][]6.8}{[1][0][]source}{}{}{}}
\makeatother
\title{Every condition admits a reservoir that stabilises the stem, by finitely many applications of strong cone avoidance}
\author{Lorenzo Carlucci, Oriola Gjetaj, Quentin Le Houérou, Ludovic Levy Patey}
\date{arxiv:2412.11598v3; CC BY 4.0}
\begin{document}
\maketitle
\noindent\textit{Editorial note.} The counted unit is the article's own Lemma lem:fs-stabilizing at source lines 1144-1147 together with its proof at source lines 1148-1150, a straight adjacent pairing in the pinned file. The pinned file is a concatenation of the authors' main and macro files; it is the article's fixed version arXiv:2412.11598v3 under CC BY 4.0.\par\medskip
\noindent\textit{Editorial note.} This source ships no separate preamble file, so the retained preamble is sliced mechanically from the pinned file by the lane and is shared, byte for byte, with the other cards of this source; it carries the authors' full macro block so that no macro of a retained span is left undefined.\par\medskip
\noindent\textit{Editorial note.} Three context spans make the section's standing material available: the convention at line 1119 that fixes the bounded-constrained colouring and its bounding function, the reminder at line 1135 that the symbol used in the statement and in the proof denotes the n-th Catalan number, and the author's definition at lines 1138-1142 of what it means for a set to stabilise a finite stem. All three are context, not counted units, and none of them overlaps the counted interval.\par\medskip
\noindent\textit{Editorial note.} The retained proof quotes one proposition of the same article, the strong cone-avoidance statement whose own environment at line 1083 carries no proof; it is declared as a non-load-bearing navigation pointer with the article's own section-wise number 6.8, which the lane counted over the numbered environments of section 6.\par\medskip
\noindent\textit{Editorial note.} The retained proof also cites the article on thin sets for reversely well-founded partial orders at source line 1149, where it invokes strong cone avoidance of the finite-dimensional free-set principle. That entry is supplied as a task-directory fragment holding it exactly as a bibtex run on the host had formatted it out of the article's own bibliography, since the bibtex backend cannot be used in this harness; the bibliography backend is therefore none.\par\medskip
\noindent\textit{Editorial note.} The statement and proof of this lemma are written as a single paragraph by the author; nothing inside them is omitted and no part of the argument points at an unreproduced statement other than the two items above.\par\medskip
\noindent\textit{Allowed standard prerequisites (editorial).}\begin{itemize}
\item Conditions, stabilisation of a stem and strong cone avoidance. The section works with conditions, each a triple of a colouring constrained by a computable bounding function, a finite stem and an infinite reservoir computing a fixed set D; a reservoir stabilises a stem when for each tuple of the stem and each length below the bounding function at the least element of the tuple the image of the reservoir on that length lands in a fixed finite set of short tuples of the stem; and the section proves separately a strong cone-avoidance statement for the finite-dimensional free-set principle, for colourings whose colours are finite sets rather than natural numbers, with the Catalan number as the bound on the number of colours. The author's own definitions and that proposition are the allowed inputs.. The counted lemma is the step that turns a condition into one whose reservoir stabilises the stem without leaving the cone of D, and it is invoked by every later step of the diagonalisation.
\end{itemize}\medskip
\ifdefined\hypertarget\hypertarget{sourcepointer-1}{}\fi
\noindent\textit{Source reference.} Original proposition 6.8: sources/logic-fl05-ramsey-like-theorems-for-the-schre/pinned-original-source.tex, source lines 993--995. Navigation pointer invoked by the retained proof for the strong cone-avoidance statement that it iterates finitely many times; the cited statement is proved elsewhere in the same article and its proof is not reproduced inside the span. The number is the article's own section-wise number for that proposition, 6.8, counted over the numbered environments of section 6; it is not a copied or guessed value.\par\medskip
In what follows, fix a coloring $f \in \P$ which is $b_f$-constrained, for some computable function $b_f : \NN \to \NN$.

In what follows, recall that $C_n$ stands for the $n$th Catalan number.

\begin{definition}
A set $X$ \emph{stabilizes} $\sigma$ if for every $t \in [\sigma]^{\leq h(\cdot)}$ with $t \neq \emptyset$ and every~$n \leq h(\min t) - |t|$,
there is a set $I_{t,n} \subseteq [\sigma]^{\leq b_f(\min t) \times C_n}$ such that for every~$u \in [X]^n$,
$f(t, u) \cap \sigma \subseteq I_{t, n}$.
\end{definition}

\begin{lemma}\label[lemma]{lem:fs-stabilizing}
For every $\sigma \in [\NN]^{<\omega}$ and every infinite set~$X$ such that $\max \sigma < \min X$ and $D \not \leq_T X$,
there is an infinite set $Y \subseteq X$ stabilizing $\sigma$ such that $D \not \leq_T Y$.
\end{lemma}

\begin{proof}
For every~$t \in [\sigma]^{\leq h(\cdot)}$ with $t \neq \emptyset$ and every~$n \leq h(\min t) - |t|$, let $g_{t, n} : [X]^n \to [\sigma]^{\leq b_f(\min t)}$ be defined by $g_{t, n}(u) = f(t \cdot u) \cap \sigma$. By finitely many successive applications of strong cone avoidance of $\RT^n_{<\infty, C_n}$ (see Cholak and Patey~\cite{cholak2020thin}), there is an infinite subset $Y \subseteq X$ such that $D \not \leq_T Y$ and for every~$t \in [\sigma]^{\leq h(\cdot)}$ with $t \neq \emptyset$ and every~$n \leq h(\min t) - |t|$, $|g_{t, n}([Y]^n)| \leq C_n$. Note that here, a color is an element of $[\sigma]^{\leq b_f(\min t)}$ instead of a natural number. Let $I_{t, n} = \bigcup g_{t, n}([Y]^n)$. Then $|I_{t, n}| \leq b_f(\min t) \times C_n$. By definition, for every~$u \in [Y]^n$, $f(t \cdot u) \cap \sigma = g_{t, n}(u) \subseteq I_{t, n}$.
\end{proof}

\begin{thebibliography}{10}

\bibitem{cholak2020thin}
P.~Cholak and L.~Patey.
\newblock Thin set theorems and cone avoidance.
\newblock {\em Trans. Amer. Math. Soc.}, 373(4):2743--2773, 2020.

\end{thebibliography}
\end{document}
